% GENERATED FILE. Edit canonical lessons, then run npm run generate:books.
% canonical-source-hash: fnv1a64:11e775adf99de0b3
% canonical-lessons: ZH-C167-qiantian, ZH-C167-fasheng, ZH-C167-meishi, ZH-W167-shi-matter, ZH-C167-tingdian, ZH-C167-shengyin

\chapter{The Day Before Yesterday}
\label{ch:zh-the-day-before-yesterday}
\hlchaptermodality{167}

\begin{chapteropening}
\textbf{By the end of this chapter:} \emph{I can say what happened the day before yesterday: a power cut, a noise, and that it was nothing.}

Complete the last lesson of chapter 167: I can say what happened the day before yesterday: a power cut, a noise, and that it was nothing.
\end{chapteropening}

\section[\texorpdfstring{\zh{前天} (qiántiān)}{前天 (qiántiān)}]{\zh{前天} (qiántiān) — the day before yesterday}

\label{lesson:ZH-C167-qiantian}



\begin{quote}
Before the new one: say the Mandarin for tea leaves, then the Mandarin for to make tea.
\end{quote}

\subsection*{You'll want to know: \zh{前天}}
\textbf{\zh{前天}} — \emph{qiántiān} — \textquotedblleft{}the day before yesterday\textquotedblright{}.

\begin{quote}
\textbf{\zh{前天是星期一}} \emph{qiántiān shì xīngqīyī} — The day before yesterday was Monday.
\end{quote}

\begin{cousinweb}
\textbf{\zh{前天}} (\emph{qiántiān}) is built from \zh{前} 'before' and \zh{天} 'day', both inherited from Middle Chinese, \emph{dzen} and \emph{then} (Guangyun, in Baxter's transcription).
\end{cousinweb}

\begin{grammarlens}[title={What you've built}]
One more word for talking about the past.
\end{grammarlens}

\subsection*{Your turn}
\begin{itemize}
\raggedright
  \item \emph{Say it:} \emph{qiántiān}
  \item \emph{Say it:} \emph{qiántiān}, once more
  \item \emph{Say it:} \emph{qiántiān shì xīngqīyī} — The day before yesterday was Monday.
  \item \emph{Recall:} say the Mandarin for tea leaves, then the Mandarin for to make tea, then say \emph{qiántiān} again
\end{itemize}

\begin{tcolorbox}[breakable,colback=teal!4,colframe=teal!35!black,title={Before you move on}]
What does \zh{前天} mean? (\textquotedblleft{}The day before yesterday\textquotedblright{}.) Say it once more.
\end{tcolorbox}
\section[fāshēng]{fāshēng — to happen; to take place}

\label{lesson:ZH-C167-fasheng}



\begin{quote}
Before the new one: say the Mandarin for to make tea, then the Mandarin for the day before yesterday.
\end{quote}

\subsection*{You'll want to know: fāshēng}
\emph{fāshēng} — \textquotedblleft{}to happen; to take place\textquotedblright{}.

\begin{quote}
\emph{qiántiān wǎnshang fāshēng le shénme} — What happened the night before last?
\end{quote}

Its characters come later; for now it is a word for the ear.

\begin{cousinweb}
\emph{fāshēng} is built from two parts inherited from Middle Chinese (Guangyun, in Baxter's transcription): \emph{pjot}, glossed 'to rise' and 'to start', the word in \emph{chūfā}, and \emph{srjaeng}, \zh{生} 'to be born': to start up and come into being.
\end{cousinweb}

\begin{grammarlens}[title={What you've built}]
One more word for talking about the past.
\end{grammarlens}

\subsection*{Your turn}
\begin{itemize}
\raggedright
  \item \emph{Say it:} \emph{fāshēng}
  \item \emph{Say it:} \emph{fāshēng}, once more
  \item \emph{Say it:} \emph{qiántiān wǎnshang fāshēng le shénme} — What happened the night before last?
  \item \emph{Recall:} say the Mandarin for to make tea, then the Mandarin for the day before yesterday, then say \emph{fāshēng} again
\end{itemize}

\begin{tcolorbox}[breakable,colback=teal!4,colframe=teal!35!black,title={Before you move on}]
What does fāshēng mean? (\textquotedblleft{}To happen; to take place\textquotedblright{}.) Say it once more.
\end{tcolorbox}
\section[\texorpdfstring{\zh{没事} (méishì)}{没事 (méishì)}]{\zh{没事} (méishì) — it's nothing; to be all right}

\label{lesson:ZH-C167-meishi}



\begin{quote}
Before the new one: say the Mandarin for the day before yesterday, then the Mandarin for to happen.
\end{quote}

\subsection*{You'll want to know: \zh{没事}}
\textbf{\zh{没事}} — \emph{méishì} — \textquotedblleft{}it's nothing; to be all right\textquotedblright{}.

\begin{quote}
\textbf{\zh{没事我很好}} \emph{méishì wǒ hěn hǎo} — It's nothing. I'm fine.
\end{quote}

The new character, in pieces: \zh{事} — eight strokes: a horizontal, a small \zh{口}, a horizontal that turns down, a long horizontal, a shorter one, and last a vertical with a hook down through them all. The next lesson puts it in your hand.

\begin{cousinweb}
\textbf{\zh{没事}} (\emph{méishì}) is built from \zh{没} 'not to have' and \zh{事} 'matter', both inherited from Middle Chinese, \emph{mwot} and \emph{dzriH} (Guangyun, in Baxter's transcription). The dictionary glosses the first 'to sink' and the second 'to employ' and 'to set up'; 'a matter' is the everyday sense, not the gloss. \emph{Mwot} would regularly give mò, not méi.
\end{cousinweb}

\begin{grammarlens}[title={What you've built}]
One more word for talking about the past.
\end{grammarlens}

\subsection*{Your turn}
\begin{itemize}
\raggedright
  \item \emph{Say it:} \emph{méishì}
  \item \emph{Say it:} \emph{méishì}, once more
  \item \emph{Say it:} \emph{méishì wǒ hěn hǎo} — It's nothing. I'm fine.
  \item \emph{Recall:} say the Mandarin for the day before yesterday, then the Mandarin for to happen, then say \emph{méishì} again
\end{itemize}

\begin{tcolorbox}[breakable,colback=teal!4,colframe=teal!35!black,title={Before you move on}]
What does \zh{没事} mean? (\textquotedblleft{}It's nothing; to be all right\textquotedblright{}.) Say it once more.
\end{tcolorbox}
\section[\texorpdfstring{\zh{事} (shì)}{事 (shì)}]{\zh{事} — a bar, \zh{口}, three bars and a long hook}

\label{lesson:ZH-W167-shi-matter}



\begin{quote}
The word first, then the character inside it.

\begin{itemize}
\raggedright
  \item \emph{Recall:} say \emph{méishì} — \textbf{R1}, one lesson back, before it decays
\end{itemize}

\begin{quote}
\zh{口} \zh{有}
\end{quote}

\emph{Read it:} each one aloud, with nothing said first — \emph{kǒu}, \emph{yǒu}
\end{quote}

\subsection*{Script — a bar, \zh{口}, three bars and a long hook}
\hlblockfigure{\detokenize{figures/ZH-W167-shi-matter-filmstrip.pdf}}{How \zh{事} is written, stroke by stroke}

\begin{quote}
\zh{事}
\end{quote}

\textbf{One shape, eight strokes, threaded on a long hook.}

\textbf{A horizontal across the top.}

\textbf{Then a small \zh{口}, three strokes,} as you write it.

\textbf{Then three bars:} a horizontal that turns down at its right end, a long horizontal across, and a shorter one under it.

\textbf{Last, a vertical with a hook,} from the top bar down through everything: the longest stroke.

1 + 3 + 3 + 1 = eight.

\subsection*{Your turn}
\begin{itemize}
\raggedright
  \item \emph{Write it:} \zh{事}, all eight, with the model in front of you
  \item \emph{Cover:} write it again from memory, and count the pieces, not the strokes
\end{itemize}

\begin{tcolorbox}[breakable,colback=teal!4,colframe=teal!35!black,title={Before you move on}]
How many strokes in \zh{事}? (\textbf{Eight}.) Which piece is second? (\textbf{A small \zh{口}}.) Which stroke comes last? (\textbf{The long vertical with the hook}.)

Source: \href{https://raw.githubusercontent.com/chanind/hanzi-writer-data/68d10a4b21150cae5e1ebbd223eed289cf32d90c/data/\%E4\%BA\%8B.json}{Hanzi Writer Data \zh{事}, pinned stroke paths and medians}.
\end{tcolorbox}
\section[tíngdiàn]{tíngdiàn — to have a power cut}

\label{lesson:ZH-C167-tingdian}



\begin{quote}
Before the new one: say the Mandarin for to happen, then the Mandarin for it's nothing.
\end{quote}

\subsection*{You'll want to know: tíngdiàn}
\emph{tíngdiàn} — \textquotedblleft{}to have a power cut\textquotedblright{}.

\begin{quote}
\emph{qiántiān wǎnshang tíngdiàn le} — There was a power cut the night before last.
\end{quote}

Its characters come later; for now it is a word for the ear.

\begin{cousinweb}
\emph{tíngdiàn} is built from \emph{tíng} 'to stop' and \zh{电} 'lightning, electricity', both inherited from Middle Chinese, \emph{deng} and \emph{denH} (Guangyun, in Baxter's transcription).
\end{cousinweb}

\begin{grammarlens}[title={What you've built}]
One more word for talking about the past.
\end{grammarlens}

\subsection*{Your turn}
\begin{itemize}
\raggedright
  \item \emph{Say it:} \emph{tíngdiàn}
  \item \emph{Say it:} \emph{tíngdiàn}, once more
  \item \emph{Say it:} \emph{qiántiān wǎnshang tíngdiàn le} — There was a power cut the night before last.
  \item \emph{Recall:} say the Mandarin for to happen, then the Mandarin for it's nothing, then say \emph{tíngdiàn} again
\end{itemize}

\begin{tcolorbox}[breakable,colback=teal!4,colframe=teal!35!black,title={Before you move on}]
What does tíngdiàn mean? (\textquotedblleft{}To have a power cut\textquotedblright{}.) Say it once more.
\end{tcolorbox}
\section[shēngyīn]{shēngyīn — sound; voice; noise}

\label{lesson:ZH-C167-shengyin}



\begin{quote}
Before the new one: say the Mandarin for it's nothing, then the Mandarin for to have a power cut.
\end{quote}

\subsection*{You'll want to know: shēngyīn}
\emph{shēngyīn} — \textquotedblleft{}sound; voice; noise\textquotedblright{}.

\begin{quote}
\emph{wàibian yǒu shénme shēngyīn} — What's that noise outside?
\end{quote}

Its characters come later; for now it is a word for the ear.

\begin{cousinweb}
\emph{shēngyīn} is built from \emph{shēng} and \zh{音}, both inherited from Middle Chinese, \emph{syeng} and \emph{'im} (Guangyun, in Baxter's transcription), which glosses the first with this very word and the second as sound made with measure.
\end{cousinweb}

\begin{grammarlens}[title={What you've built}]
One more word for talking about the past.
\end{grammarlens}

\subsection*{Your turn}
\begin{itemize}
\raggedright
  \item \emph{Say it:} \emph{shēngyīn}
  \item \emph{Say it:} \emph{shēngyīn}, once more
  \item \emph{Say it:} \emph{wàibian yǒu shénme shēngyīn} — What's that noise outside?
  \item \emph{Recall:} say the Mandarin for it's nothing, then the Mandarin for to have a power cut, then say \emph{shēngyīn} again
\end{itemize}

\begin{tcolorbox}[breakable,colback=teal!4,colframe=teal!35!black,title={Before you move on}]
What does shēngyīn mean? (\textquotedblleft{}Sound; voice; noise\textquotedblright{}.) Say it once more.
\end{tcolorbox}
